% ============================================================
%  SECCIÓN — CÁLCULOS GEOTÉCNICOS (ARQUETA)
% ============================================================
\section{Informe de cálculos geotécnicos}
\label{sec:arqueta_geotecnia}

La arqueta de fangos mezcla tiene dimensiones interiores
$L\times B\times H = 4{,}00\times3{,}00\times3{,}50$~m,
con paredes, losa de fondo y tapadera de espesor uniforme $e = 0{,}30$~m
(HA-30/B/20/IIb+Qc). Las dimensiones exteriores resultan:
\[
  L_\text{ext} = 4{,}60~\text{m},\quad
  B_\text{ext} = 3{,}60~\text{m},\quad
  A_\text{cim} = 4{,}60\times3{,}60 = 16{,}56~\text{m}^2
\]

% -----------------------------------------------------------
\subsection{Datos de partida y cargas sobre la cimentación}
\label{subsec:cargas_arqueta}
% -----------------------------------------------------------

\subsubsection{Volúmenes y pesos de hormigón}

Las paredes frontales ($\times2$) tienen longitud $L_\text{ext}=4{,}60$~m;
las laterales encajan entre ellas con longitud
$B_\text{ext}-2e = 3{,}00$~m. La altura de las paredes es $H=3{,}50$~m (interior),
pues la losa y la tapadera se contabilizan por separado:

\begin{align*}
  V_\text{pared frente/fondo} &= 4{,}60\times3{,}50\times0{,}30 = 4{,}830~\text{m}^3
    \quad(\times2)\\
  V_\text{pared lat.}         &= 3{,}00\times3{,}50\times0{,}30 = 3{,}150~\text{m}^3
    \quad(\times2)\\
  V_\text{losa} = V_\text{tapa} &= 4{,}60\times3{,}60\times0{,}30 = 4{,}968~\text{m}^3
    \quad(\times2)\\
  V_{HA} &= 2\times4{,}830 + 2\times3{,}150 + 2\times4{,}968
          = \mathbf{25{,}896}~\text{m}^3
\end{align*}
\[
  W_{HA} = V_{HA}\times\gamma_{HA} = 25{,}896\times2{,}50
         = \resultado{64{,}74~\text{Tn}}
\]

\subsubsection{Peso del fango (depósito lleno)}

\[
  W_\text{fango} = \gamma_f\times L\times B\times H
    = 1{,}10\times4{,}00\times3{,}00\times3{,}50
    = \resultado{46{,}20~\text{Tn}}
\]

\subsubsection{Carga total y presión de contacto}

\begin{align*}
  W_\text{total} &= W_{HA} + W_\text{fango}
    = 64{,}74 + 46{,}20 = \resultado{110{,}94~\text{Tn}}\\[4pt]
  \sigma_\text{contacto} &= \frac{W_\text{total}}{A_\text{cim}}
    = \frac{110{,}94}{16{,}56} = \resultado{6{,}70~\text{Tn/m}^2}
\end{align*}

\begin{table}[H]
  \centering
  \caption{Resumen de cargas transmitidas al terreno --- Arqueta de fangos mezcla}
  \label{tab:cargas_arqueta}
  \begin{tabular}{@{}lrcc@{}}
    \toprule
    \textbf{Componente} & \textbf{V (m³)} & $\boldsymbol{\gamma}$ \textbf{(Tn/m³)} & \textbf{Peso (Tn)}\\
    \midrule
    Pared de frente ($4{,}60\times3{,}50\times0{,}30$) & 4,830 & 2,50 & 12,08\\
    Pared de fondo  ($4{,}60\times3{,}50\times0{,}30$) & 4,830 & 2,50 & 12,08\\
    Pared lat.\ izq.\ ($3{,}00\times3{,}50\times0{,}30$) & 3,150 & 2,50 & 7,88\\
    Pared lat.\ der.\ ($3{,}00\times3{,}50\times0{,}30$) & 3,150 & 2,50 & 7,88\\
    Losa de fondo ($4{,}60\times3{,}60\times0{,}30$)   & 4,968 & 2,50 & 12,42\\
    Tapadera ($4{,}60\times3{,}60\times0{,}30$)        & 4,968 & 2,50 & 12,42\\
    \midrule
    \textbf{Total hormigón} & \textbf{25,896} & 2,50 & \textbf{64,74}\\
    Fango interior ($4{,}00\times3{,}00\times3{,}50$)  & 42,000 & 1,10 & 46,20\\
    \midrule
    \textbf{TOTAL} & --- & --- & \textbf{110,94}\\
    \bottomrule
  \end{tabular}
\end{table}

% -----------------------------------------------------------
\subsection{Selección de la cota de cimentación}
\label{subsec:seleccion_cota_arqueta}
% -----------------------------------------------------------

\subsubsection*{Análisis del perfil estratigráfico}

Se evalúa la aptitud de cada estrato como plano de apoyo:

\begin{itemize}
  \item \textbf{E-1a (0 a $-2{,}50$~m) — Arena limosa no saturada, $N_\text{SPT}=9$,
    $c'=0$, $\phi'=29°$, $E=6{,}894$~MPa.}
    Para $\sigma_\text{contacto}=6{,}70$~Tn/m², el estrato es \textbf{admisible}:
    capacidad portante drenada holgada, asientos cuantificables y excavación en seco
    (por encima del N.F.\ a $-2{,}50$~m).
  \item \textbf{E-1b ($-2{,}50$ a $-3{,}50$~m) — Arena limosa saturada.}
    Bajo el N.F. Cimentar aquí exigiría agotamiento continuo con riesgo de
    sifonamiento y arrastres.
  \item \textbf{E-2 ($-3{,}50$ a $-8{,}00$~m) — Arcillas versicolores, $C_u=1{,}00$~kg/cm².}
    Alta capacidad portante no drenada, pero exige $D_f\geq4{,}0$~m y atravesar la
    capa saturada E-1b, con entibación y agotamiento.
  \item \textbf{E-3 ($-8{,}00$ a $-25{,}00$~m) — Arcillas margosas firmes.}
    Solo accesible mediante pilotes.
\end{itemize}

\subsubsection*{Hipótesis $D_f=4{,}0$~m (apoyo en E-2) — descartada}

El empotramiento mínimo recomendado por CTE DB SE-C ($\geq 0{,}5$–$1{,}0$~m
dentro del estrato) lleva a $D_f\geq4{,}0$~m. Inconvenientes detectados:
\begin{itemize}
  \item Excavación $1{,}8$~m bajo el N.F.: agotamiento continuo durante toda la fase.
  \item Atravesar E-1b saturado: riesgo de sifonamiento, arrastres y desestabilización
        de taludes; requiere entibación o tablestacas provisionales.
  \item Sin beneficio técnico: cimentar en E-1a proporciona $FS_{CP}>15$ y
        asientos $<3$~cm.
\end{itemize}

\begin{notabox}[Hipótesis Df = 4{,}0 m en E-2 --- descartada]
  Cimentar en E-2 obliga a atravesar el N.F.\ en zona medioambientalmente
  protegida, con entibación y agotamiento continuados, sin mejora significativa
  de las verificaciones geotécnicas.
\end{notabox}

\subsubsection*{Cota adoptada: $D_f=1{,}80$~m (apoyo en E-1a)}

Se sigue la convención del CTE DB SE-C: $D_f$ es la profundidad del
\emph{intradós} (cara inferior) de la losa, medida desde rasante.
Con $e_\text{losa}=0{,}30$~m, el trasdós superior queda a cota $-1{,}50$~m.

\begin{supuestobox}[Cota de cimentación adoptada]
  \begin{itemize}
    \item Estrato de apoyo: \textbf{E-1a} (arena limosa no saturada).
    \item Cota intradós / base de excavación: $-1{,}80$~m
          $\Rightarrow$ \textbf{$D_f=1{,}80$~m}.
    \item Cota trasdós superior de la losa: $-1{,}50$~m.
    \item Excavación en \textbf{seco}: la base queda $0{,}70$~m por encima
          del N.F.\ ($-2{,}50$~m).
    \item Coronación de la arqueta:
          $-1{,}50 + 3{,}50 + 0{,}30 = \mathbf{+2{,}30}$~m sobre rasante.
  \end{itemize}
\end{supuestobox}

% -----------------------------------------------------------
\subsection{Estados límite aplicables}
\label{subsec:estados_limite_arqueta}
% -----------------------------------------------------------

\begin{itemize}
  \item \textbf{Hundimiento} — SÍ aplica. Verificado en \S\ref{subsec:cap_portante_arqueta}
        mediante Brinch-Hansen drenado.
  \item \textbf{Asientos} — SÍ aplican. Calculados en \S\ref{subsec:asientos_arqueta}
        mediante Steinbrenner y Boussinesq.
  \item \textbf{Flotación} — SE VERIFICA en \S\ref{subsec:flotacion}.
        El N.F.\ ($-2{,}50$~m) queda por debajo del intradós ($-1{,}80$~m);
        en servicio normal no existe subpresión.
  \item \textbf{Vuelco} — NO aplica. Caja simétrica; los empujes activos de
        paredes opuestas se anulan dos a dos.
  \item \textbf{Deslizamiento} — NO aplica. La resultante horizontal neta es nula.
  \item \textbf{Sifonamiento} — NO aplica. La arqueta en servicio es impermeable;
        no existe gradiente hidráulico.
\end{itemize}

% -----------------------------------------------------------
\subsection{Empujes del terreno sobre las paredes}
\label{subsec:empujes_arqueta}
% -----------------------------------------------------------

Las paredes están en contacto con E-1a ($\phi'=29°$, $c'=0$,
$\gamma=1{,}82$~Tn/m³). La altura de pared enterrada es
$H_\text{ent}=1{,}50$~m (de rasante al trasdós superior de la losa).
Dado que el N.F.\ ($-2{,}50$~m) está por debajo del fondo de excavación
($-1{,}80$~m), \textbf{no existe presión hidrostática lateral} sobre las paredes.

\paragraph{Coeficiente de empuje activo:}
\[
  K_a = \frac{1-\sin\phi'}{1+\sin\phi'}
      = \frac{1-0{,}4848}{1+0{,}4848} = \frac{0{,}5152}{1{,}4848}
      = \resultado{0{,}347}
\]

\paragraph{Presión activa en la base de la pared enterrada:}
\[
  \sigma'_{h,a} = K_a\,\gamma_{1a}\,H_\text{ent}
    = 0{,}347\times1{,}82\times1{,}50 = \resultado{0{,}947~\text{Tn/m}^2}
\]

\paragraph{Empuje activo total por metro lineal de pared:}
\[
  E_a = \tfrac{1}{2}\,K_a\,\gamma_{1a}\,H_\text{ent}^2
      = \tfrac{1}{2}\times0{,}347\times1{,}82\times1{,}50^2
      = \resultado{0{,}710~\text{Tn/ml}}
\]
aplicado a $H_\text{ent}/3=0{,}50$~m sobre la base de la pared.
Dato de entrada para el dimensionamiento estructural de las paredes.

% -----------------------------------------------------------
\subsection{Capacidad portante --- Brinch-Hansen (drenado)}
\label{subsec:cap_portante_arqueta}
% -----------------------------------------------------------

Se aplica la fórmula de Brinch-Hansen en condiciones \textbf{drenadas}
(largo plazo) sobre la losa $B^*\times L^*=3{,}60\times4{,}60$~m.
Con $c'=0$, el término cohesivo se anula: $q_u = T_2 + T_3$.

\paragraph{Factores de capacidad de carga ($\phi'=29°$):}
\begin{align*}
  N_q &= e^{\pi\tan\phi'}\cdot\tan^2\!\left(45°+\tfrac{\phi'}{2}\right)
       = 5{,}704\times2{,}882 = \resultado{16{,}44}\\
  N_c &= \frac{N_q-1}{\tan\phi'} = \frac{15{,}44}{0{,}5543} = 27{,}86\\
  N_\gamma &= 1{,}5\,(N_q-1)\tan\phi' = 1{,}5\times15{,}44\times0{,}5543
           = \resultado{12{,}84}
\end{align*}

\paragraph{Factores de forma ($B^*/L^*=3{,}60/4{,}60=0{,}7826$):}
\begin{align*}
  s_q      &= 1+\frac{B^*}{L^*}\sin\phi'
            = 1+0{,}7826\times0{,}4848 = \resultado{1{,}379}\\
  s_\gamma &= 1-0{,}4\,\frac{B^*}{L^*}
            = 1-0{,}4\times0{,}7826 = \resultado{0{,}687}
\end{align*}

\paragraph{Factores de profundidad ($D_f/B^*=1{,}80/3{,}60=0{,}500$):}
\begin{align*}
  d_q      &= 1+2\tan\phi'(1-\sin\phi')^2\,\frac{D_f}{B^*}
            = 1+2\times0{,}5543\times0{,}2654\times0{,}500
            = \resultado{1{,}147}\\
  d_\gamma &= 1{,}0
\end{align*}

\paragraph{Sobrecarga al nivel del intradós de la losa:}
\[
  q_0 = \gamma_{1a}\times D_f = 1{,}82\times1{,}80 = \resultado{3{,}28~\text{Tn/m}^2}
\]

\paragraph{Peso específico ponderado $\gamma_k$ bajo el intradós (zona $B^*=3{,}60$~m):}

El N.F.\ se sitúa a $0{,}70$~m del intradós, por lo que la zona de influencia
abarca tres tramos con densidades distintas:

\begin{table}[H]
  \centering
  \begin{tabular}{@{}lcccc@{}}
    \toprule
    \textbf{Tramo bajo intradós} & \textbf{Cotas} & \textbf{e (m)} & \textbf{Estrato} & $\boldsymbol{\gamma}$ (Tn/m³)\\
    \midrule
    0,00–0,70~m & $-1{,}80$ a $-2{,}50$ & 0,70 & E-1a (no sat.) & 1,82\\
    0,70–1,70~m & $-2{,}50$ a $-3{,}50$ & 1,00 & E-1b ($\gamma'=0{,}98$) & 0,98\\
    1,70–3,60~m & $-3{,}50$ a $-5{,}40$ & 1,90 & E-2 ($\gamma'=0{,}78$)  & 0,78\\
    \bottomrule
  \end{tabular}
\end{table}

\[
  \gamma_k = \frac{0{,}70\times1{,}82 + 1{,}00\times0{,}98 + 1{,}90\times0{,}78}{3{,}60}
           = \frac{1{,}274+0{,}980+1{,}482}{3{,}60}
           = \frac{3{,}736}{3{,}60} = \resultado{1{,}038~\text{Tn/m}^3}
\]

\paragraph{Carga de hundimiento:}
\begin{align*}
  T_1 &= 0 \quad (c'=0)\\
  T_2 &= q_0\,N_q\,s_q\,d_q
       = 3{,}28\times16{,}44\times1{,}379\times1{,}147
       = \resultado{85{,}29~\text{Tn/m}^2}\\
  T_3 &= \tfrac{1}{2}\,\gamma_k\,B^*\,N_\gamma\,s_\gamma\,d_\gamma
       = 0{,}5\times1{,}038\times3{,}60\times12{,}84\times0{,}687\times1{,}0
       = \resultado{16{,}48~\text{Tn/m}^2}\\[4pt]
  q_u          &= T_2+T_3 = 85{,}29+16{,}48 = \resultado{101{,}77~\text{Tn/m}^2}\\
  \sigma_\text{adm} &= q_u/3 = 101{,}77/3 = \resultado{33{,}92~\text{Tn/m}^2}
\end{align*}

\paragraph{Verificación:}
\begin{equation}
  FS_\text{real} = \frac{q_u}{\sigma_\text{contacto}}
    = \frac{101{,}77}{6{,}70}
    = \resultado{15{,}2 \gg 3{,}0 \quad \checkmark}
  \label{eq:fs_cp_arqueta}
\end{equation}

\begin{resultadobox}[Capacidad portante --- Arqueta]
  El estrato~E-1a ofrece $q_u=101{,}77$~Tn/m² y $\sigma_\text{adm}=33{,}92$~Tn/m².
  Con $\sigma_\text{contacto}=6{,}70$~Tn/m² el factor de seguridad es $FS=15{,}2\gg3{,}0$.
  No existe riesgo de hundimiento.
\end{resultadobox}

% -----------------------------------------------------------
\subsection{Verificación de flotación}
\label{subsec:flotacion}
% -----------------------------------------------------------

El N.F.\ del enunciado ($-2{,}50$~m) queda \textbf{$0{,}70$~m por debajo}
del intradós de la losa ($-1{,}80$~m). \textbf{No existe columna de agua bajo
la cimentación en condiciones de servicio normales.}

\[
  V_\text{sumergido} = 0~\text{m}^3
  \;\Rightarrow\;
  E_\text{flot} = \gamma_w\times V_\text{sumergido} = \mathbf{0~\text{Tn}}
\]

\begin{resultadobox}[Verificación de flotación]
  $E_\text{flot}=0 < W_{HA}=64{,}74$~Tn $\;\Rightarrow\;$
  \textbf{Sin riesgo de flotación} \checkmark.
  No se requieren anclajes ni medidas adicionales.
\end{resultadobox}

\begin{notabox}[Hipótesis conservadora — N.F.\ en rasante]
  Si el N.F.\ ascendiese eventualmente hasta rasante (hipótesis pésima), el
  empuje de Arquímedes sería
  $F_\text{Arq} = 1{,}00\times(4{,}60\times3{,}60\times1{,}80) = 29{,}81$~Tn,
  frente a $W_{HA}=64{,}74$~Tn.
  El factor de seguridad resultante sería
  $FS_\text{flot}=64{,}74/29{,}81=\mathbf{2{,}17}>1{,}05$~\checkmark.
  Con el drenaje perimetral proyectado y la clase medioambiental de la zona,
  este escenario se considera improbable.
\end{notabox}

% -----------------------------------------------------------
\subsection{Presión neta y efecto de compensación}
\label{subsec:presion_neta}
% -----------------------------------------------------------

La presión neta es la diferencia entre la carga aplicada y la sobrecarga
geostática previa en el plano de cimentación; es la magnitud que genera
nuevos asientos respecto al estado natural del terreno:

\[
  \sigma_\text{neta} = \sigma_\text{contacto} - q_0
    = 6{,}70 - 3{,}28 = \resultado{3{,}42~\text{Tn/m}^2}
\]

La compensación alcanza el $49$\,\%: la excavación libera casi la mitad del
peso aportado por la estructura y el fango, reduciendo considerablemente
los asientos.

% -----------------------------------------------------------
\subsection{Cálculo de asientos}
\label{subsec:asientos_arqueta}
% -----------------------------------------------------------

Los asientos se calculan con la \textbf{presión neta}
$\sigma_\text{neta}=3{,}42$~Tn/m², que representa el incremento real de carga
sobre el estado geostático preexistente.

\[
  S_\text{total} = S_i + S_c
\]

\subsubsection{Asiento inmediato --- Steinbrenner (4 sub-rectángulos)}

Se aplica el método de la esquina: la losa se divide en cuatro sub-rectángulos de
$B'\times L'=1{,}80\times2{,}30$~m y el asiento en el centro se obtiene multiplicando
por~4. Con $n=L'/B'=2{,}30/1{,}80=1{,}278$, $E_s=703{,}0$~Tn/m² y $\nu=0{,}25$:

\paragraph{Factor de influencia $F_1$:}
\begin{align*}
  F_1 &= \frac{2}{\pi}\left[n\ln\!\left(\frac{\sqrt{n^2+1}+1}{n}\right)
         + \ln\!\left(n+\sqrt{n^2+1}\right)\right]\\
      &= \frac{2}{\pi}\left[1{,}278\times\ln(2{,}052) + \ln(2{,}900)\right]
       = \frac{2}{\pi}\times1{,}984 = \resultado{1{,}263}
\end{align*}

\paragraph{Asiento en esquina del sub-rectángulo:}
\[
  S_\text{esq} = \frac{\sigma_\text{neta}\cdot B'}{2\,E_s}(1-\nu^2)\,F_1
    = \frac{3{,}42\times1{,}80}{2\times703{,}0}\times0{,}9375\times1{,}263
    = \resultado{0{,}518~\text{cm}}
\]

\paragraph{Asiento en el centro:}
\[
  S_i = 4\times S_\text{esq} = 4\times0{,}518 = \resultado{2{,}07~\text{cm}}
\]

\subsubsection{Asiento de consolidación --- Estrato E-2}

La influencia se extiende hasta el Estrato~2 (arcillas versicolores).
Profundidad al centro de E-2 desde el intradós de la losa:
$z = 5{,}75 - 1{,}80 = 3{,}95$~m.

Con $m=L'/z=2{,}30/3{,}95=0{,}582$ y $n=B'/z=1{,}80/3{,}95=0{,}456$
se obtiene $I_B=0{,}0870$ (factor de Boussinesq):

\[
  \Delta\sigma_z = \sigma_\text{neta}\times4\times I_B
    = 3{,}42\times4\times0{,}0870 = 1{,}19~\text{Tn/m}^2
\]
\[
  S_c = \frac{\Delta\sigma_z\cdot H_{E_2}}{E_{E_2}}
      = \frac{1{,}19\times4{,}50}{1223{,}6} = \resultado{0{,}44~\text{cm}}
\]

donde $H_{E_2}=4{,}50$~m (de $-3{,}50$ a $-8{,}00$~m) y
$E_{E_2}=12{,}000$~kPa~$=1223{,}6$~Tn/m².

\subsubsection{Asiento total y verificación}

\begin{equation}
  S_\text{total} = S_i + S_c = 2{,}07 + 0{,}44
    = \resultado{2{,}51~\text{cm} \leq 5{,}00~\text{cm} \quad \checkmark}
  \label{eq:asiento_total_arqueta}
\end{equation}

\begin{resultadobox}[Asientos --- Arqueta]
  El asiento total estimado ($2{,}51$~cm) cumple holgadamente el límite
  orientativo de $5{,}0$~cm para estructuras de hormigón.
  La compensación del $49$\,\% sobre la carga bruta reduce a la mitad los
  asientos respecto de un análisis sin descuento geostático.
\end{resultadobox}

% -----------------------------------------------------------
\subsection{Alternativa de cimentación profunda --- Pilotes}
\label{subsec:pilotes_arqueta}
% -----------------------------------------------------------

Aunque la cimentación superficial satisface todas las verificaciones, se
analiza como alternativa la \textbf{cimentación por pilotes hormigonados
in situ con lodos bentoníticos} (CPI-8, CTE DB SE-C sección~5.1.2), con punta
empotrada en el Estrato~E-3.

\subsubsection{Geometría del pilote adoptado}

$D=0{,}40$~m, $L=10{,}0$~m, cabeza a rasante, punta a $-10{,}0$~m
(empotramiento $2{,}0$~m en E-3, $>2D=0{,}80$~m \checkmark):
\[
  A_p = \pi D^2/4 = 0{,}1257~\text{m}^2;\quad
  p   = \pi D = 1{,}257~\text{m}
\]

Tramos activos: E-2 ($-3{,}50$ a $-8{,}00$~m), $L_2=4{,}50$~m,
$C_u=10$~Tn/m²; E-3 ($-8{,}00$ a $-10{,}00$~m), $L_3=2{,}00$~m,
$C_u=15{,}2$~Tn/m². Se desprecia fuste en E-1 (arena).

\subsubsection{Resistencia a corto plazo — no drenado}

\paragraph{Punta:}
\[
  q_p = N_p\,C_{u,\text{punta}} = 9\times15{,}2 = 136{,}8~\text{Tn/m}^2
\]

\paragraph{Fuste} ($\tau_f = 100\,C_u/(100+C_u)$~kPa):
\begin{align*}
  \tau_{f,E_2} &= 50{,}0~\text{kPa} = 5{,}10~\text{Tn/m}^2\\
  \tau_{f,E_3} &= 60{,}3~\text{kPa} = 6{,}15~\text{Tn/m}^2
\end{align*}

\paragraph{Resistencias características y de cálculo:}
\begin{align*}
  R_{pk}  &= 136{,}8\times0{,}1257 = 17{,}19~\text{Tn}\\
  R_{fk}  &= 5{,}10\times1{,}257\times4{,}50 + 6{,}15\times1{,}257\times2{,}00
           = 28{,}85+15{,}46 = 44{,}30~\text{Tn}\\
  R_{ck}  &= 17{,}19+44{,}30 = \resultado{61{,}50~\text{Tn/pilote}}\\
  R_{cd}  &= R_{ck}/\gamma_R = 61{,}50/3{,}0 = \resultado{20{,}50~\text{Tn/pilote}}
\end{align*}
($\gamma_R=3{,}0$ para pilotes perforados con lodos, Anejo~F.2 CTE.)
Gobierna sobre el análisis a largo plazo ($R_{cd,LP}=36{,}10$~Tn/pilote).

\subsubsection{Número de pilotes y verificación del grupo}

\[
  n_\text{mín} = \frac{W_\text{total}}{R_{cd}}
    = \frac{110{,}94}{20{,}50} = 5{,}41
    \;\Rightarrow\; \textbf{n=6~pilotes} \quad(2\times3~\text{en planta})
\]

Separaciones en encepado $4{,}60\times3{,}60$~m con recubrimiento $1{,}5D=0{,}60$~m:
\[
  \text{sep}_B = (3{,}60-1{,}20)/1 = 2{,}40~\text{m} > 3D=1{,}20~\text{m}\;\checkmark;\quad
  \text{sep}_L = (4{,}60-1{,}20)/2 = 1{,}70~\text{m} > 3D\;\checkmark
\]

Sin efecto de grupo ($\eta=1{,}0$):
\begin{align*}
  FS_\text{grupo}    &= \frac{\eta\,n\,R_{ck}}{W_\text{total}}
    = \frac{1{,}0\times6\times61{,}50}{110{,}94}
    = \resultado{3{,}33 \geq 3{,}0 \quad \checkmark}\\
  N_\text{pilote} &= 110{,}94/6 = 18{,}49~\text{Tn}
    \leq R_{cd}=20{,}50~\text{Tn}\;\checkmark
\end{align*}

Asiento estimado del grupo (Vesic, $B_g=\text{sep}_B+D=2{,}80$~m):
\[
  S_g = \sqrt{B_g/D}\times S_{i,\text{pilote}}
      = \sqrt{2{,}80/0{,}40}\times3{,}3~\text{mm}
      \approx \mathbf{8{,}7~\text{mm}} < 25~\text{mm}\;\checkmark
\]

\subsubsection{Comparativa y descarte de la solución profunda}

\begin{table}[H]
  \centering
  \caption{Comparativa cimentación superficial vs.\ profunda --- Arqueta}
  \label{tab:comp_cimentacion_arqueta}
  \begin{tabular}{@{}lll@{}}
    \toprule
    \textbf{Criterio} & \textbf{Superficial} & \textbf{Profunda (pilotes)}\\
    \midrule
    $FS$ hundimiento          & 15,2 (drenado)     & 3,33 (grupo, ND)\\
    Asiento total             & $\approx$25~mm      & $\approx$9~mm\\
    Flotación                 & Sin riesgo          & Sin riesgo (anclaje)\\
    Profundidad excavación    & 1,80~m              & 10,0~m (perforación)\\
    Zona saturada atravesada  & No                  & Sí ($-2{,}50$ a $-10{,}0$~m)\\
    Complejidad ejecución     & Baja                & Alta (lodos bentoníticos)\\
    Afección al N.F.           & Nula                & Reducida\\
    Coste relativo            & $1\times$ (ref.)    & $2{,}5$–$4\times$\\
    Plazo de ejecución        & 2–3 semanas         & 4–6 semanas\\
    \bottomrule
  \end{tabular}
\end{table}

\begin{notabox}[Decisión: cimentación superficial adoptada]
  La alternativa de pilotes cumple técnicamente ($FS_\text{grupo}=3{,}33$,
  $S_g=8{,}7$~mm), pero supone un sobrecoste de $2{,}5$–$4\times$ y mayor
  plazo y complejidad sin que exista ninguna insuficiencia en la solución
  superficial. Se descarta la cimentación profunda.
\end{notabox}

% -----------------------------------------------------------
\subsection{Síntesis de verificaciones geotécnicas}
\label{subsec:sintesis_geotecnia_arqueta}
% -----------------------------------------------------------

\begin{table}[H]
  \centering
  \caption{Síntesis final de verificaciones geotécnicas --- Arqueta de fangos mezcla}
  \label{tab:sintesis_geotecnia_arqueta}
  \begin{tabular}{@{}llll@{}}
    \toprule
    \textbf{Verificación} & \textbf{Valor calculado} & \textbf{Límite} & \textbf{Resultado}\\
    \midrule
    $\sigma_\text{contacto}$ bruta & 6,70~Tn/m² & $<\sigma_\text{adm}$ & ---\\
    $\sigma_\text{adm}$ drenada    & 33,92~Tn/m² & — & ---\\
    \textbf{$FS$ hundimiento (drenado)} & \textbf{15,2} & $\geq3{,}0$ & \textbf{OK \checkmark}\\
    $\sigma_\text{neta}$ & 3,42~Tn/m² & compensación 49\% & ---\\
    \textbf{Asiento total estimado} & \textbf{25,1~mm} & $\leq50$~mm & \textbf{OK \checkmark}\\
    \textbf{Flotación} (N.F.\ a $-2{,}50$~m) & \textbf{Sin subpresión} & — & \textbf{OK \checkmark}\\
    Empuje activo por m.l. & 0,710~Tn/ml & — & dato estructural\\
    Tipo de cimentación & Losa superficial & $D_f=1{,}80$~m, E-1a & ---\\
    \bottomrule
  \end{tabular}
\end{table}

La cimentación superficial mediante losa HA-30/B/20/IIb+Qc, apoyada en
E-1a a $D_f=1{,}80$~m, satisface holgadamente todas las verificaciones
del CTE DB SE-C. Se descarta la alternativa de cimentación profunda por
suponer un sobrecoste y sobreplazo no justificados técnicamente.
